\subsection{Smoothing and its joint complex Hessian}
\label{sec:chern-envelope-smoothing}
The metric-Laplacian estimate in a big class will come from a quadratic bound for the smoothing error. Ordinary coordinate convolution does not define a global family. We construct a kernel from finite holomorphic jets of the exponential map and prove the joint estimate in the base point and the radius that will be needed to correct its small negative part.

For the calculations in this and the next two subsections put
\[
 \mathcal D=\mathrm i\partial\bar\partial=2\pi\ddc,
 \qquad \alpha=2\pi\theta,\qquad \Omega=2\pi\omega.
\]
Thus $\alpha+\mathcal Du\geq0$ is precisely the original $\theta$-psh condition. Multiplying forms by $2\pi$ only changes the constants in the estimates. Let $h$ be the coefficient metric of $\Omega$, and assume $n\geq1$ during the smoothing calculation. Square brackets with a superscript $2$ denote evaluation of the associated Hermitian form on a vector; thus $\Gamma[v,\eta]^2$ is its value on the product tangent vector $(v,\eta)$. Brackets applied to a Hermitian matrix denote contraction with that matrix.

\begin{lemma}\label{lem:chern-joint-smoothing}
For some $r_0>0$ there is a positive averaging kernel $\Psi_u(x,w)$, smooth for $0<|w|<r_0$, invariant under rotations of $w$, and linear in $u\in\PSH(X,\theta)$, with the following properties. Constants may depend on the fixed geometry, $\sup_Xu$ and $\|u\|_{L^1}$, but not on the radius. For a fixed sufficiently large $K$, put
\[
 H_u(x,r)=\Psi_u(x,r)+Kr^2,\qquad
 \lambda_u(x,r)=r\partial_rH_u(x,r).
\]
The function $H_u(x,\mathrm e^s)$ is strictly convex and increasing in $s$, tends to the canonical $u(x)$ as $s\to-\infty$, and $\lambda_u\geq0$. At a point where $u$ is locally bounded, $\lambda_u(x,r)\to0$.

At each $x$ use normal coordinates of uniformly fixed size, let $p=u+q$ be a local psh potential with $\mathcal Dq=\alpha$, and let $\mu_p$ be its positive Euclidean trace measure. Put $d=2n$ and
\[
 m_u(x,r)=1+r^{-d}\mu_p(B(x,4r)).
\]
All estimates below hold uniformly over the fixed normal charts. For a base tangent vector $v$ and a radius tangent vector $\eta\in\mathbb C$, the joint Levi form satisfies
\begin{align}\label{eq:chern-joint-mass-estimate}
 (\alpha+\mathcal D_{x,w}H_u)[v,\eta]^2
 \geq{}&-\bigl(A\lambda_u+C(r^2+r^3m_u)\bigr)|v|_h^2\notag\\
       &-C(r+r^2m_u)|v|_h|\eta|-Cr^2|\eta|^2,
       \qquad r=|w|.
\end{align}
In particular $m_u\leq Cr^{-2}$ gives
\begin{equation}\label{eq:chern-initial-joint-bound}
 (\alpha+\mathcal D_{x,w}H_u)[v,\eta]^2
 \geq-(A\lambda_u+Cr)|v|_h^2-C|v|_h|\eta|-C|\eta|^2.
\end{equation}
If $m_u(x,r)\leq M(x)/r$ with $M(x)\geq1$, then the stronger estimate is
\begin{equation}\label{eq:chern-improved-joint-bound}
 (\alpha+\mathcal D_{x,w}H_u)[v,\eta]^2
 \geq-(A\lambda_u+CM(x)r^2)|v|_h^2
                    -CM(x)r|v|_h|\eta|-CM(x)r|\eta|^2.
\end{equation}
Finally, on a compact set where $u$ is bounded, $\Psi_u(x,r)$ differs from ordinary radial convolution in these normal coordinates by $O(r^2)$. For a $C^1$ function $f$ with bounded weak complex Hessian,
\begin{equation}\label{eq:chern-regular-obstacle-error}
 |\Psi_f(x,r)-f(x)|\leq C_f r^2,
\end{equation}
uniformly on $X$, with $C_f$ controlled by its value, gradient and complex-Hessian bounds.
For every locally integrable $u$, the infinitesimal generator is
\begin{align}\label{eq:chern-smoothing-generator}
 \frac{\Psi_u(\cdot,r)-u}{r^2}\longrightarrow
 a_\chi D_hu\quad\hbox{in distributions},&\qquad
 D_hu=\sum h^{i\bar j}u_{i\bar j},\notag\\
 &
 a_\chi=\frac1n\int|\zeta|^2\chi(|\zeta|^2)\,\mathrm d\zeta>0.
\end{align}
Here $\chi$ is the normalized radial weight in the kernel constructed below.
\end{lemma}
\begin{proof}
We give the kernel and its derivative calculation. Fix the jet order $N=8$. Normal holomorphic coordinates and a holomorphic frame can be chosen so that, at the chosen center,
\[
 \langle e_l,e_m\rangle_h
   =\delta_{lm}-\sum_{j,k}c_{jklm}z_j\bar z_k+O(|z|^3).
\]
For a K\"ahler metric the frame can be the coordinate frame. The quadratic and cubic holomorphic changes of coordinates remove the first metric derivatives and the pure holomorphic second jets: their coefficients are symmetric because $h$ is the complex Hessian of a local potential. These changes have coefficients bounded by fixed metric derivatives on a finite atlas, so the chart sizes and the bounds for the remaining finite jets are uniform. The $c_{jklm}$ are the curvature coefficients, as follows by differentiating $\Gamma^m_{jl}=h^{m\bar a}\partial_jh_{l\bar a}$.

\Needspace{5\baselineskip}
Let $\exp_x(\zeta)$ be the time-one map of the geodesic equation
\[
 \ddot z^m+\Gamma^m_{jl}(z)\dot z^j\dot z^l=0.
\]
Its existence, smoothness and finite Taylor expansion on a uniformly small neighborhood follow from its integral equation: iteration on a short time interval is a contraction for bounded smooth coefficients, and differentiation of the integral equation gives its successive parameter derivatives. Subdivide the time interval to reach time one. Keep only the holomorphic fiber terms of its Taylor expansion through degree $N$. These jets are compatible under holomorphic coordinate changes, since substitution by a holomorphic target map and a complex linear fiber map preserves the holomorphic part of a jet.

They have a global smooth representative $F(x,\zeta)$. To see this without an infinite-jet extension theorem, take the local polynomial representatives $F_\nu$ on a finite cover. Write each in the ordinary exponential chart at $x$ as $\exp_x(a_\nu(x,\zeta))$, average $a_\nu$ with a smooth base partition, and apply $\exp_x$ again. The $F_\nu$ differ by $O(|\zeta|^{N+1})$ on overlaps, with the corresponding differentiated bounds, so this construction preserves all jets through degree $N$. Also $F(x,0)=x$ and $D_\zeta F(x,0)$ is the identity. The inverse function theorem and compactness give a uniformly small radius on which every fiber map is a real diffeomorphism. In the normal coordinates at its center,
\begin{equation}\label{eq:chern-normal-kernel-map}
 F(x,\zeta)=x+\zeta+O(|\zeta|^3),
 \qquad D_\zeta F=I+O(|\zeta|^2).
\end{equation}
Indeed the only nonholomorphic cubic term of the geodesic expansion is removed, and its holomorphic quadratic and cubic terms vanish in these coordinates.

Choose
\[
 \chi(t)=\frac{C_n}{(1-t)^2}\mathrm e^{-1/(1-t)}\quad(t<1),
 \qquad \chi(t)=0\quad(t\geq1),
\]
with $\int_{\mathbb C^n}\chi(|\zeta|^2)\,\mathrm d\zeta=1$, and let
$\chi_1(t)=-C_n\mathrm e^{-1/(1-t)}$ for $t<1$, zero otherwise.
Thus $\chi_1'=\chi$ and $-\chi_1/\chi=(1-t)^2$. Define
\[
 \Psi_u(x,w)=\int_{T_xX}u(F(x,w\zeta))
                       \chi(|\zeta|_h^2)\,\mathrm d\zeta_h.
\]
It is rotation invariant by changing $\zeta$ to $\mathrm e^{\mathrm it}\zeta$. The fiber change of variables to $y=F(x,w\zeta)$ writes it as integration against a smooth compactly supported kernel for $w\ne0$. It is consequently smooth even when $u$ has poles.

Here is the joint Hessian calculation. We use multi-indices $I,J$ with $2\leq|I|,|J|\leq N$ and write $1_l$ for the $l$th coordinate multi-index. A factor $I_ls^{I-1_l}$ is zero when $I_l=0$; the same convention applies to its conjugate and to $J$. A local factorization of the finite jets has the form
\[
 F(z,\xi)=a(z,\rho(z,\xi)),\qquad
 \rho_m=\xi_m+\sum_I\left(\sum_kd_{Ikm}\bar z_k
                   +\sum_{j,k}e_{Ijkm}z_j\bar z_k\right)\xi^I+R_m.
\]
Here $a$ is holomorphic, with linear part $z+\xi$; its inverse in $\xi$ exists near zero. To obtain the factorization, put all purely holomorphic terms through base degree two and fiber degree $N$ in $a$, and apply its local inverse to $F$. The displayed terms are exactly those of the resulting jet contributing to its first base derivatives or its mixed base second derivative at zero. The other terms in $R_m$ are pure quadratic base terms, terms of base degree at least three, or the higher fiber remainder. All their relevant differentiated bounds follow from Taylor's formula, and all displayed coefficients are uniformly bounded.

Set $s=w^{-1}\rho(z,w\xi)$. In these variables
\[
 \Psi_u(z,w)=\int u(a(z,ws))\chi(A(z,w,s))B(z,w,s)\,\mathrm ds.
\]
Writing only the terms that contribute to the joint Levi jet at $z=0$, the squared norm $A$ and the density $B$ are
\begin{align*}
 A={}&|s|^2-\sum c_{jklm}z_j\bar z_k s_l\bar s_m
       -2\operatorname{Re}\sum d_{Ikm}w^{|I|-1}s^I\bar s_m\bar z_k\\
    &-2\operatorname{Re}\sum e_{Ijkm}w^{|I|-1}s^I\bar s_mz_j\bar z_k
       +\sum d_{Ikm}\overline d_{Jjm}
         w^{|I|-1}\bar w^{|J|-1}s^I\bar s^Jz_j\bar z_k+R_A,\\
 B={}&1-\sum c_{jkll}z_j\bar z_k
       -2\operatorname{Re}\sum d_{Ikm}w^{|I|-1}I_ms^{I-1_m}\bar z_k\\
    &-2\operatorname{Re}\sum e_{Ijkm}w^{|I|-1}I_ms^{I-1_m}z_j\bar z_k\\
    &+\sum d_{Ikm}\overline d_{Jjl}
       w^{|I|-1}\bar w^{|J|-1}I_mJ_ls^{I-1_m}\bar s^{J-1_l}z_j\bar z_k+R_B.
\end{align*}
These identities follow by substituting $\xi=s-\sum d_Iw^{|I|-1}s^I\bar z+\cdots$ in the norm and its volume density. The inverse substitution's pure quadratic antiholomorphic terms and their conjugates have zero relevant Levi jet and are included in the remainders. After any two relevant derivatives, the higher fiber remainders contribute $O(r^{N-2})$ against a bounded-support kernel.

For completeness, the integration by parts eliminating first derivatives of $u$ uses the identity
\[
 \sum_m\partial_{s_m}
      \left(\chi_1(|s|^2)\frac{I_m}{|I|}\bar s^{I-1_m}\right)
       =\chi(|s|^2)\bar s^I.
\]
\Needspace{18\baselineskip}
Equivalently one may differentiate once more in $s_l$ and use
$\sum_lI_ls^{I-1_l}s_l=|I|s^I$ to group the weight derivatives. Applying these product rules to the displayed $A,B$ and integrating twice, with no boundary term because the weights vanish at $|s|=1$, gives the positive-square identity below. Its coefficients can be recorded explicitly. Put
\begin{align*}
 \Xi^m={}&\sum_{I,j,l}\frac{I_m}{|I|}\overline d_{Ijl}
                 \bar w^{|I|-2}\bar s_l\bar s^{I-1_m}v_j,\\
 U_{ml}={}&-\frac{\chi_1}{\chi}\left(
      \sum_{j,k}c_{jklm}v_j\bar v_k
      +2\sum_{I,j,k}e_{Ijkm}w^{|I|-1}
                         \frac{I_l}{|I|}s^{I-1_l}v_j\bar v_k\right.\\
 &\hspace{29mm}\left.
      +2\sum_{I,k}d_{Ikm}(|I|-1)w^{|I|-2}
                          \frac{I_l}{|I|}s^{I-1_l}\eta\bar v_k\right)\\
 &+\sum_{I,J,j,k}d_{Ikm}\overline d_{Jjl}
                    w^{|I|-1}\bar w^{|J|-1}s^I\bar s^Jv_j\bar v_k.
\end{align*}
The quotient is evaluated at $|s|^2$. Replace $U$ by its Hermitian part. If $P=(p_{l\bar m})\geq0$ is the complex Hessian of the local psh potential $p$, evaluated at $a(0,ws)$, the identity is
\begin{equation}\label{eq:chern-kernel-square-identity}
 \mathcal D\Psi_p[v,\eta]^2
   =\int P\bigl[Da(v,\eta s+r^2\Xi)\bigr]^2\chi\,\mathrm ds
       +r^2\int P\bigl[Da(U-r^2\Xi\Xi^*)\bigr]\chi\,\mathrm ds+E.
\end{equation}
In the second integral the differential acts on the vertical Hermitian matrix. One can check the grouping directly: differentiation of $u(a(z,ws))\chi(A)B$ gives its positive Hessian term, two first-derivative terms from $\partial(\chi(A)B)$, and its value times $\partial\bar\partial(\chi(A)B)$. The two first-derivative terms, after the preceding product rules, are exactly the cross terms of the first square. The curvature, $e$, mixed $d$ and quadratic $d\bar d$ terms are exactly the four terms displayed in $U$. Completing the square subtracts $r^2\Xi\Xi^*$ from that matrix. This verifies the identity using all the differentiated weight terms.

We spell out the bounds for its remainders and coefficients. The local Lelong--Jensen formula in \cref{thm:euclidean-jensen-lelong}, together with finite trace mass on a fixed larger chart, bounds
\[
 \mu_p(B(x,t))\leq Ct^{d-2},\qquad
   t^{-d}\int_{B(x,t)}|u|\leq C(1+|\log t|).
\]
For the second estimate integrate the bounded normalized trace mass in the spherical-mean derivative; a fixed outer annulus has a lower average controlled by $\|u\|_{L^1}$, and the upper bound for $u$ bounds its positive part. These arguments are uniform on the finite atlas. Thus $m_u\leq Cr^{-2}$. The twice differentiated higher-jet kernel remainders are bounded by $Cr^{N-2}(1+|\log r|)$. Replacing $a$ by $F$ and the associated variables changes the current kernels by $O(r^N)$; positivity bounds this error by $Cr^Nm_u$. This comparison is a change of variables in their integrals, not a pointwise comparison of Hessians at different points. With $N=8$, both errors are bounded by $Cr^2(|v|^2+|\eta|^2)$.

The first integral of \eqref{eq:chern-kernel-square-identity} is nonnegative. The coefficients above give $|\Xi|\leq C|v|$, while the leading part of $U$ is its curvature matrix multiplied by $-\chi_1/\chi$. The remaining terms of $U$ have bounds $C(r|v|^2+|v||\eta|+r^2|v|^2)$. Their integrals, and the square correction, are consequently bounded below by
\[
 -A\lambda_0(x,r)|v|^2
       -C(r^3m_u+r^2)|v|^2-Cr^2m_u|v||\eta|-Cr^2|\eta|^2,
\]
where
\[
 \lambda_0(x,r)=2r^2\int(-\chi_1(|s|^2))
                           \operatorname{tr}P(x+rs)\,\mathrm ds.
\]
The bounded curvature norm gives $A$; changes of the differential $Da$ and of variables cost at most $Cr^4m_u\leq Cr^2$. Subtracting the smoothing of the fixed smooth local potential $q$ adds horizontal $O(r^2)$ and mixed $O(r)$ terms and a bounded vertical term. Choose $K$ once to absorb the latter with a strictly positive vertical remainder. For $v=0$, the positive-square formula shows directly that $\partial_w\partial_{\bar w}\Psi_u\geq-C$; thus this choice also makes $H_u(x,\mathrm e^s)$ strictly convex in $s$.

We finally relate $\lambda_0$ to the radial derivative of our actual kernel. The normal-coordinate change \eqref{eq:chern-normal-kernel-map} changes the normalized value kernel from Euclidean convolution by $O(r^2)$. Its value error is $O(r^2(1+|\log r|))$. For the vertical Hessian use \eqref{eq:chern-kernel-square-identity} with $v=0$: the positive current kernels differ from the Euclidean vertical kernel by $O(r^2)$, hence their integral error is $O(1)$ by $m_u\leq Cr^{-2}$. Smooth $q$ and the added $Kr^2$ give bounded errors as well. The radial identity
\[
 \partial_w\partial_{\bar w}a(r)
                       =\frac1{4r}\partial_r(r\partial_ra)
\]
then shows that the difference of the radial derivatives is $O(r^2)$ up to a constant of integration. The value error tends to zero, so that constant must vanish: a nonzero constant would give a logarithmic value error. Euclidean integration by parts gives its radial derivative as $\lambda_0$ for $p$, and $q$ changes it by $O(r^2)$. Therefore $|\lambda_u-\lambda_0|\leq Cr^2$. This proves \eqref{eq:chern-joint-mass-estimate} and both of its stated consequences.

The upper bound for $u$ bounds $H_u$ above as $s\to-\infty$. Its convexity consequently makes its slope nonnegative. The Euclidean convolution tends to the canonical $u$, and the value error tends to zero, proving the asserted pointwise limit. Where $u$ is locally bounded that limit is finite; a nonnegative increasing convex slope must then tend to zero at $-\infty$.

On a bounded neighborhood the same value-kernel comparison has error $O(r^2\|u\|_\infty)$, proving the final local assertion. If $f$ has bounded complex Hessian, the ordinary radial convolution error is $O(r^2)$: its Euclidean Laplacian is bounded, and the spherical-mean identity integrated against $\chi$ gives the bound. The displacement error in \eqref{eq:chern-normal-kernel-map} is $O(r^3\|\nabla f\|_\infty)$. This proves \eqref{eq:chern-regular-obstacle-error}.

For \eqref{eq:chern-smoothing-generator}, work in a fixed coordinate chart. The holomorphic fiber jets of $F$ have no mixed quadratic term. The first moments and the pure holomorphic second moments of the radial weight vanish, whereas
$\int\zeta_i\bar\zeta_j\chi=a_\chi\delta_{ij}$. Taylor expansion therefore gives
\[
 \Psi_g=g+a_\chi r^2\sum h^{i\bar j}g_{i\bar j}+O(r^3)
\]
for smooth $g$, with bounds for every fixed number of base derivatives on smaller charts. To obtain the distributional statement, apply the same expansion to the adjoint kernel. Choose a smooth local orthonormal frame and, for each $\zeta$ in the support of $\chi$, change the base variable by $x\mapsto F(x,r\zeta)$. Its differential is $I+O(r)$, so it is a diffeomorphism on a slightly smaller chart for small $r$. Taylor expansion of its inverse and Jacobian, followed by integration in $\zeta$, gives, with respect to coordinate Lebesgue measure,
\[
 \Psi_r^*\varphi=\varphi+
 a_\chi r^2\sum_{i,j}\partial_i\partial_{\bar j}
                  (h^{i\bar j}\varphi)+O(r^3).
\]
The displayed coefficients can also be identified by pairing with an arbitrary smooth $g$ in the preceding expansion; the Taylor remainder for the inverse and its Jacobian is uniform in every fixed $C^k$ norm. Pairing this adjoint formula with locally integrable $u$ proves \eqref{eq:chern-smoothing-generator}. This keeps track of the moving metric in the kernel.

All calculations for nonsmooth psh $p$ are justified by local decreasing convolution: the current integrals have smooth compactly supported kernels, positivity gives the uniform bounds above, and their distributional limits give exactly the displayed identities and inequalities.
\end{proof}

\subsection{A minimum with a varying upper radius}
\label{sec:varying-radius-minimum}
The regularization used in a big class must become narrower near its analytic pole set. A minimum over these varying radii needs a boundary argument in addition to the product version of Kiselman's principle. The following form records the curvature error of the upper radius explicitly.

\begin{lemma}\label{lem:varying-log-radius-minimum}
Let $D$ be a coordinate domain with a fixed smooth positive form $\Omega$, let $b$ be smooth and real on $D$, and suppose $-C_b\Omega\leq\mathrm i\partial\bar\partial b\leq C_b\Omega$. Let $H(z,s)$ be smooth for $s$ below, and in a neighborhood of, $b(z)$, with $H_s\geq0$ and $H_{ss}>0$. Regard $H$ as a function of $(z,\zeta)$ with $s=\operatorname{Re}\zeta$. Assume $H$ is bounded below locally uniformly in $z$ as $s\to-\infty$.

Fix $c>0$ and a smooth real $(1,1)$-form $\Lambda$ on $D$, and let $\pi_D$ denote projection to the base $D$. Suppose that at every point with $s\leq b(z)$ and $H_s\leq c$, the joint Levi form satisfies
\[
 \mathrm i\partial\bar\partial_{z,\zeta}H\geq-\pi_D^*\Lambda.
\]
Then the finite constrained minimum
\[
 v(z)=\inf_{s\leq b(z)}\bigl(H(z,s)-cs\bigr)
\]
is continuous and satisfies, as distributions,
\[
 \mathrm i\partial\bar\partial v\geq-\Lambda-cC_b\Omega.
\]
The infimum includes the upper endpoint. In particular a constant upper radius has no additional curvature error.
\end{lemma}
\begin{proof}
For $\varepsilon>0$ minimize
\[
 H(z,s)-cs-\varepsilon\log(b(z)-s),\qquad s<b(z).
\]
This expression tends to infinity at both ends of the fiber. Near $-\infty$ the positive linear term $-cs$ dominates the logarithm, by the lower bound for $H$. Its second $s$ derivative is $H_{ss}+\varepsilon/(b-s)^2>0$, so it has a unique minimizer $s_\varepsilon(z)$. The implicit function theorem makes this minimizer, and its minimum $v_\varepsilon$, smooth. Its critical equation is
\[
 H_s-c+\frac\varepsilon{b-s}=0.
\]
At that point $0\leq H_s\leq c$ and $0\leq\sigma:=\varepsilon/(b-s)\leq c$. The joint Levi form of the logarithmic barrier equals a nonnegative square minus $\sigma\,\mathrm i\partial\bar\partial b$. Hence the joint Levi form of the minimized expression, at its minimizer, is bounded below by $-\pi_D^*(\Lambda+cC_b\Omega)$.

The Schur-complement computation in the proof of \cref{thm:kiselman-invariant-exhaustion} applies at this point without requiring a Levi bound elsewhere in the fiber. In detail the joint matrix is
$\bigl(\begin{smallmatrix}C&B/2\\B^*/2&A/4\end{smallmatrix}\bigr)$ with $A>0$, and differentiating the critical equation gives the minimum matrix $C-BA^{-1}B^*$. Adding the stated horizontal lower-bound form makes the full matrix nonnegative; evaluation on $(q,-2A^{-1}B^*q)$ gives the same lower bound for the minimum. Thus
\[
 \mathrm i\partial\bar\partial v_\varepsilon
                                  \geq-\Lambda-cC_b\Omega.
\]

We justify the limiting minimum as well. On a compact base neighborhood, $b$ is bounded and evaluation at $s=b-1$ gives a uniform upper bound for $v_\varepsilon$ when $0<\varepsilon\leq1$. The locally uniform lower bound for $H$ and linear growth of $-cs-\log(b-s)$ give a uniform lower bound for all $s_\varepsilon$. The unpenalized minima have the same lower bound. Both problems may therefore be restricted to a common compact interval below their varying upper endpoint. The unpenalized value is continuous by uniform continuity on that compact set: writing $s=b(z)-t$ makes the upper endpoint fixed, and the range of $t$ can be bounded uniformly.

On this range the penalty is bounded below by a quantity tending uniformly to zero. For an upper bound evaluate at an unpenalized minimizer if it is at least $\sqrt\varepsilon$ below $b$, and otherwise at $s=b-\sqrt\varepsilon$. The first evaluation costs at most $\varepsilon|\log\sqrt\varepsilon|$; in the second the unpenalized value changes by $o(1)$ uniformly, by its smoothness on the compact range. Consequently $v_\varepsilon\to v$ locally uniformly. Passing the Levi inequalities to distributions proves the assertion.
\end{proof}

\subsection{Big-class envelopes on the ample locus}
\label{sec:big-envelope-metric-regularity}
The loss of positivity of $\theta$ is confined by a strict potential with analytic singularities. Near its pole set we will shrink the smoothing radius. The first estimate for that radius comes from the defining maximality of the envelope; high powers of the local divisor equations then turn the variable coefficients in the joint Hessian estimate into uniform ones.

\Needspace{17\baselineskip}
\begin{theorem}\label{thm:big-metric-laplacian-rooftop}
Suppose $\{\theta\}$ is big. Let $f_1,\ldots,f_r$ be continuous real functions with bounded distributional Laplacians for the fixed metric $\omega$, where $r\geq1$ is an integer. Then
\[
 u=P_\theta\bigl(\min_j f_j\bigr)
\]
has bounded metric Laplacian on every compact subset of the ample locus
\[
 \operatorname{Amp}(\{\theta\})=X\setminus E_{\mathrm{nK}}(\{\theta\}).
\]
Its complex Hessian has locally bounded coefficients there, and
$u\in W^{2,p}_{\mathrm{loc}}$ for every $1\leq p<\infty$, hence
$u\in C^{1,\gamma}_{\mathrm{loc}}$ for every $0<\gamma<1$.
\end{theorem}
\begin{proof}
In complex dimension zero the statement concerns functions on finitely many points and is immediate. Assume $n\geq1$.

First reduce the obstacles to one with bounded complex Hessian. Choose a K\"ahler form $\kappa=L\omega$ with $\kappa\geq\theta$, and put
$f=P_\kappa(\min_jf_j)$. By \cref{thm:kahler-metric-laplacian-rooftop}, $f$ is $C^1$ and has bounded complex Hessian. Every $\theta$-psh function below $\min_jf_j$ is $\kappa$-psh, and is therefore below $f$. Conversely $f\leq\min_jf_j$. These two inclusions of competitor sets give $P_\theta(f)=u$.

Choose a strict potential with gentle analytic singularities and pole set $A=E_{\mathrm{nK}}(\{\theta\})$, using \cref{thm:positivity-loci,lem:logarithmic-cutoff-sum}. Principalize its ideal by the smooth blowups of \cref{thm:analytic-resolution-input}, obtaining $\pi:Y\to X$, an isomorphism over $X\setminus A$. The pulled-back potential has a smooth divisorial remainder. The proof of \cref{prop:kahler-smooth-center-blowup}, applied successively, supplies
\[
 \omega_Y=\pi^*\omega+\sum_\ell\varepsilon_\ell\eta_\ell>0,
\]
where $\eta_\ell$ is the curvature of $\mathcal O(-E_\ell)$ and the effective exceptional divisors and positive coefficients can be rational. If the pulled-back strict current is at least $\delta\pi^*\omega$, add $\delta\sum_\ell\varepsilon_\ell\log|s_{E_\ell}|^2$ to its potential, taking $\delta>0$ rational below its original margin. Since
\[
 \ddc\log|s_{E_\ell}|^2=[E_\ell]+\eta_\ell,
\]
the resulting current is at least $\delta\omega_Y$. It still has a smooth divisorial remainder and positive rational divisor coefficients. Rescale the measuring K\"ahler form to make this lower bound equal to that form.

The potential bijection in \cref{prop:smooth-target-psh-modification} identifies the obstacle envelopes upstairs and downstairs. The pullback of $f$ is $C^1$ with bounded complex Hessian: the chain rule bounds its coefficients off the exceptional divisor, and a local positive potential obtained by adding a smooth background is bounded. Apply \cref{thm:cln-weighted} to that bounded potential and a local logarithmic pole defining the divisor. The logarithm is integrable for its Hessian trace, so that trace gives the divisor zero mass. Thus the bounded coefficient currents extend across it with no additional part. We can consequently work on $Y$ with the pulled-back smooth form and obstacle, and return through the isomorphism at the end.

For this proof rename $Y$ as $X$ and its data as $\theta,\omega,f$, and again put $\alpha=2\pi\theta$, $\Omega=2\pi\omega$. Add a constant to the obstacle so that $0\leq f\leq B$, with $B\geq1$, and subtract a constant from the strict reference $\phi$ so that
\begin{equation}\label{eq:big-envelope-reference-bounds}
 \alpha+\mathcal D\phi\geq\Omega,\qquad
 \phi\leq-B-1,\qquad \phi\leq u\leq f\leq B.
\end{equation}
On a finite divisor atlas, with smaller charts covering $X$, write
\[
 \phi=\sum_j a_j\log|g_j|^2+b,
 \qquad a_j>0,\quad |g_j|<\tfrac12,
\]
where $b$ is smooth. Write $D$ for the divisor and $U=X\setminus D$. In particular $\mathcal D\phi$ is uniformly bounded on $U$, although $\phi$ tends to $-\infty$ near $D$. All bounds below are uniform on this finite atlas; only explicitly displayed functions of $\phi$ are allowed to grow.

\emph{A first-order smoothing bound.}
Let $H=H_u$ and $\lambda=\lambda_u$ be the corrected smoothing of \cref{lem:chern-joint-smoothing}. Its initial joint estimate improves after adding $K_1r$. Indeed the vertical Hessian of $K_1|w|$ is $K_1/(4r)$; Young's inequality absorbs the bounded mixed term at a horizontal cost $Cr$. For $K_1$ fixed sufficiently large and $r$ small,
\begin{equation}\label{eq:big-first-corrected-joint}
 \alpha+\mathcal D(H+K_1r)\geq-(A\lambda+Cr)\Omega
\end{equation}
as a joint form; the right side is pulled back from $X$.
Fix $0<\delta<\delta_0$ and $c>0$, and define on $U$
\[
 \psi_{c,\delta}(x)=\inf_{0<t\leq\delta}
 \{H(x,t)+K_1t-K_1\delta-c\log(t/\delta)\}.
\]
The fiber minimum is finite and attained, since $H\geq u$ and $u$ is locally bounded on $U$. At a minimizing radius, including the upper endpoint, its increasing part has logarithmic slope at most $c$, so $\lambda\leq c$. Apply \cref{lem:varying-log-radius-minimum} with constant ceiling $\log\delta$ to \eqref{eq:big-first-corrected-joint}. Evaluation at $t=\delta$ and \eqref{eq:chern-regular-obstacle-error} give
\[
 \alpha+\mathcal D\psi_{c,\delta}\geq-(Ac+C\delta)\Omega,
 \qquad \psi_{c,\delta}\leq f+C\delta^2.
\]
Put $e=Ac+C\delta$ and $\mu=e/(1+e)$. Then
$(1-\mu)(\psi_{c,\delta}-C\delta^2)+\mu\phi$ is $\alpha$-psh on $U$ and is below $f$. Its bounded-above extension across $D$ exists by analytic removal, and remains below the continuous obstacle. Envelope maximality implies
\begin{equation}\label{eq:big-first-envelope-comparison}
 \psi_{c,\delta}\leq u+C\delta^2
                         +\mu(B-\phi+C\delta^2).
\end{equation}
For this last rearrangement use $\psi_{c,\delta}\leq B+C\delta^2$; thus no unbounded factor is divided by $1-\mu$.

Choose $c=\delta$ and let $t_0$ be a minimizing radius. Since $H(t_0)\geq u$ and $t_0\leq\delta$, \eqref{eq:big-first-envelope-comparison} gives, for $R=C(1+B+|\phi|)$ enlarged if necessary,
\[
 t_0\geq\delta\mathrm e^{-R(x)},\qquad
 H(x,t_0)+K_1t_0-u(x)\leq R(x)\delta.
\]
Monotonicity of $H$ and the choice $\delta=r\mathrm e^{R(x)}$ consequently prove
\begin{equation}\label{eq:big-first-smoothing-error}
 0\leq H(x,r)-u(x)\leq R(x)\mathrm e^{R(x)}r
 \quad\hbox{if }0<r<\delta_0\mathrm e^{-R(x)}.
\end{equation}
Since $\phi\leq-B-1$, we may take $R(x)=B_0|\phi(x)|$ with $B_0$ fixed large, while decreasing $\delta_0$.

This also improves the local trace-mass bound. Choose $B_0$ larger if needed. The divisor formula then implies
$\mathrm e^{-R(x)}\leq C|g_j(x)|^2$ for each local divisor equation. Bounded first derivatives of $g_j$ show that, for $|y-x|\leq S\mathrm e^{-R(x)}$ and a fixed small $S$, the ratios $|g_j(y)|/|g_j(x)|$ lie between $1/2$ and $3/2$. Hence $|\phi(y)-\phi(x)|\leq C$ there. On a chart with no divisor component, smoothness of $\phi$ gives the same oscillation bound after decreasing $S$. By \eqref{eq:big-envelope-reference-bounds}, $|u(y)|\leq C(1+|\phi(x)|)$ on that ball. The normal-coordinate kernel comparison in \cref{lem:chern-joint-smoothing} transfers \eqref{eq:big-first-smoothing-error} to the ordinary radial mean of the local psh potential $p=u+q$, with error $Cr^2(1+|\phi(x)|)$.

For clarity, the needed consequence of Jensen's formula has a positive constant independent of $x,r$:
\[
 \int p(x+r\zeta)\chi(|\zeta|^2)\,\mathrm d\zeta-p(x)
        \geq c_0 r^{2-d}\mu_p(B(x,c_1r)),
 \qquad d=2n,\quad 0<c_1<1.
\]
Indeed the Green weight for each spherical mean is positive on a smaller ball; integrate the radii in a fixed subinterval of $(c_1,1)$ where $\chi>0$. For $d=2$ the logarithmic Green weight gives the same statement with $r^{2-d}=1$. Apply this at a fixed larger multiple of $r$, decrease $S$, and use \eqref{eq:big-first-smoothing-error}. We obtain
\begin{equation}\label{eq:big-improved-trace-mass}
 m_u(x,r)\leq\frac{M(x)}r\quad(0<r<S\mathrm e^{-R(x)}),
 \qquad 1\leq M(x)\leq C R(x)\mathrm e^{CR(x)}.
\end{equation}
Thus the improved joint estimate \eqref{eq:chern-improved-joint-bound} is available on these radii.

\emph{Uniform joint estimates in divisor coordinates.}
On each chart put $G=(\prod_jg_j)^m$, with an empty product equal to one. Take the integer $m$ so large that $M|G|$ and $M|\mathrm dG|$ are uniformly bounded: the divisor formula for $R$ and \eqref{eq:big-improved-trace-mass} bound $M$ by a fixed product of negative powers of $|g_j|$, and differentiation of $G$ loses only one power of one factor. A sufficiently large positive power dominates this product, including the factor $R$. Next choose $k\geq1$ so large that
\[
 \frac{\mathrm e^{-kR(x)}}{|G(x)|}\leq T
\]
with $T$ uniform on the finite cover. This follows again from the same divisor formula. Put $\rho_\delta(x)=\delta\mathrm e^{-kR(x)}$, taking $\delta_0<S$. In the holomorphic coordinates $w=G(x)\xi$, the range $0<|w|\leq\rho_\delta(x)$ has $|\xi|\leq T\delta$.

The differential of this change sends $(v,\eta)$ to
$(v,G\eta+\xi\,\mathrm dG(v))$. Substitute it in \eqref{eq:chern-improved-joint-bound}. The bounds just obtained, $|G|\leq1$ and a uniform bound for $|\mathrm dG|$, give
\begin{align}\label{eq:big-divisor-transformed-joint}
 (\alpha+\mathcal D H(x,|G\xi|))[v,\eta]^2
 \geq{}&-(A\lambda+C_1|\xi|^2)|v|^2\notag\\
        &-C_1|G||\xi||v||\eta|
                      -C_1|G|^2|\xi||\eta|^2.
\end{align}
For example its purely horizontal error is bounded by a constant times
$M(|G|^2+|G||\mathrm dG|+|G||\mathrm dG|^2)|\xi|^2$;
its mixed coefficient is bounded by $C|G||\xi|$, and its vertical coefficient by $C|G|^2|\xi|$. These estimates use only the two bounded products $M|G|$ and $M|\mathrm dG|$.

Add $K_2|G\xi|^2$. Its Levi form is the square
$K_2|G\eta+\xi\,\mathrm dG(v)|^2$. The inequality
$|a+b|^2\geq\tfrac12|a|^2-|b|^2$ leaves a positive multiple of $K_2|G|^2|\eta|^2$, at a horizontal cost $K_2|\mathrm dG|^2|\xi|^2$. For $\delta_0$ small and $K_2$ fixed large, this absorbs the last term of \eqref{eq:big-divisor-transformed-joint}. Young's inequality absorbs its mixed term at cost $C|\xi|^2|v|^2$. We have proved the uniform joint bound
\begin{equation}\label{eq:big-second-corrected-joint}
 \alpha+\mathcal D(H(x,|G\xi|)+K_2|G\xi|^2)
                         \geq-(A\lambda+C_2\delta^2)\Omega.
\end{equation}
All constants here are independent of the divisor chart and the point of $U$.

\emph{A second-order smoothing bound.}
For $c>0$ set
\[
 \psi(x)=\inf_{0<t\leq\rho_\delta(x)}
               \{H(x,t)+K_2t^2-c\log t\}.
\]
In the coordinate $s=\log|\xi|$, its upper endpoint is
$b(x)=\log\delta-kR(x)-\log|G(x)|$.
The last summand is pluriharmonic on $U$ and
$\mathcal D R=-B_0\mathcal D\phi$ is uniformly bounded there. Thus $|\mathcal D b|\leq C_b\Omega$ with a uniform constant. The logarithmic slope of the increasing part is $\lambda+2K_2t^2$, so at the minima it is at most $c$. Applying \cref{lem:varying-log-radius-minimum} to \eqref{eq:big-second-corrected-joint}, and noting that $-c\log|G|$ has zero Levi form, proves
\begin{equation}\label{eq:big-second-minimum-curvature}
 \alpha+\mathcal D\psi\geq-C_0\Omega,
 \qquad C_0=A'c+C_2\delta^2.
\end{equation}
This application includes a minimum at the upper endpoint of the permitted radius.

Let $a=ckB_0$ and $C_3=C_f+K+K_2$, with $C_f$ from \eqref{eq:chern-regular-obstacle-error}. Evaluation at $t=\rho_\delta(x)$ gives
\[
 Q:=\psi+c\log\delta+a\phi-C_3\delta^2\leq f,
 \qquad (1+a)\alpha+\mathcal DQ\geq-C_0\Omega.
\]
The second inequality uses \eqref{eq:big-envelope-reference-bounds} in addition to \eqref{eq:big-second-minimum-curvature}. Put
\[
 \nu=\frac{C_0}{1+a+C_0},\qquad
 w=(1-\nu)Q+\nu(1+a)\phi.
\]
This is $(1+a)\alpha$-psh. It is below $f$, because $Q\leq f$, $\phi\leq0$ and $f\geq0$. Analytic removal extends it over $D$, preserving its upper bound. Consequently $w/(1+a)$ is a competitor for $u$, and $(1+a)u\geq w$. Since $u,Q\leq B$, we can rearrange this without losing an unbounded positive term:
\begin{equation}\label{eq:big-second-envelope-comparison}
 u\geq Q+\nu(1+a)\phi-(\nu+a)B.
\end{equation}

Choose $c=\delta^2$, so that $a$ and $\nu$ are $O(\delta^2)$, and let $t_0$ be a minimizing radius for $\psi(x)$. Substitute its value in \eqref{eq:big-second-envelope-comparison}. Since $H(t_0)\geq u$, and $\log(t_0/\delta)\leq0$, it yields
\begin{equation}\label{eq:big-second-radius-and-value}
 t_0\geq\delta\mathrm e^{-R_2(x)},\qquad
 0\leq H(x,t_0)+K_2t_0^2-u(x)\leq R_2(x)\delta^2,
 \qquad R_2=C(1+B+|\phi|).
\end{equation}
Indeed the nonnegative terms on the left after rearrangement are
$H(t_0)+K_2t_0^2-u$ and $-\delta^2\log(t_0/\delta)$; all terms on the right are bounded by $C\delta^2(1+B+|\phi|)$.
\Needspace{10\baselineskip}
Fix a compact neighborhood $K\Subset U$ and $M_2\geq\sup_KR_2$. For every sufficiently small $r$, choose $\delta=r\mathrm e^{M_2}$. Then $r\leq t_0$ at every $x\in K$, and monotonicity gives the uniform quadratic estimate
\begin{equation}\label{eq:big-quadratic-smoothing-error}
 0\leq H(x,r)-u(x)\leq M_2\mathrm e^{2M_2}r^2,
 \qquad x\in K.
\end{equation}

Finally divide \eqref{eq:big-quadratic-smoothing-error} by $r^2$ and use \eqref{eq:chern-smoothing-generator}. The distribution
$a_\chi D_hu+K$ lies between zero and a locally bounded constant, hence is an $L^\infty_{\mathrm{loc}}$ function by \cref{lem:dominated-measure-density}. Positivity of $\alpha+\mathcal Du$ and its bounded $h$-trace bound each diagonal measure in a fixed chart by Lebesgue measure. The matrix Cauchy--Schwarz inequality bounds its off-diagonal measures as well. Thus all its complex coefficients, and those of $\mathcal Du$, are locally bounded.

To obtain the claimed real Sobolev regularity without assuming weak first derivatives, the Euclidean Laplacian of $u$ on a smaller chart is now bounded. Multiply this bounded right side by a compactly supported cutoff, take its fundamental-solution potential, and subtract it from $u$ on the region where the cutoff equals one. The difference is distributionally harmonic, hence smooth by Weyl's lemma. The Hessian multiplier estimate in \cref{thm:hessian-multiplier-lp} puts the potential in $W^{2,p}_{\mathrm{loc}}$ for every $1<p<\infty$; $p=1$ follows on relatively compact sets by finite volume. Morrey's inequality gives $C^{1,\gamma}_{\mathrm{loc}}$ for every $\gamma<1$. The continuous representative agrees with the canonical psh one by its local submean characterization. Transferring these statements through the isomorphism $Y\setminus\pi^{-1}(A)\to X\setminus A$ proves the theorem.
\end{proof}
\index{envelope!big-class regularity}
\index{ample locus!envelope regularity}
